% =====================================================================================
\chapter{SYSTEM OVERVIEW AND CALIBRATION}
\label{chp:system}
% Material: README "Big picture", "How the pieces connect", §1 (frames), §16 (depth),
% §14 (calibration story); notes/pen_marking_plan.md; thesis_notes.md (calibration
% entries); config/pen_tool.json; notes/realtime_fp.md (the GPU split).
% =====================================================================================

\section{Hardware}
\label{sec:sys_hardware}
\begin{table}[htbp]
\caption{Hardware of the cell.}
\label{tab:hardware}
\centering\small
\begin{tabular}{@{}lll@{}}
\toprule
Component & Model & Role \\
\midrule
Manipulator & Universal Robots UR5e (factory-calibrated kinematics) & motion, F/T sensing \\
Camera & Intel RealSense D435i, eye-in-hand, off-axis & RGB-D, 1280$\times$720 \\
Tool & pen in a holder (CAD envelope) & marks instead of a torch \\
Workpieces & 8 mm MDF plates; 3D-printed curved parts & \\
Desktop & GPU (RTX 5070 Ti) & registration (SAM2, PPF, FoundationPose) \\
Laptop & GPU (RTX 4060, 8 GB) & ROS 2, tracking, planning \\
\bottomrule
\end{tabular}
\end{table}
\figplaceholder{fig:tool}{The pen tool: holder, pen, camera arm; the collision envelope
(RViz, /tool\_model/markers)}{The pen tool and its collision envelope.}

\section{Software Architecture}
\label{sec:sys_software}
\TODO{ROS 2 Jazzy nodes and what feeds what; the desktop registers, the laptop tracks
and plans; the operator's gestures.}
\SRC{README ``How the pieces connect''; notes/pipeline\_map.png; realtime\_fp.md status.}
\figplaceholder{fig:architecture}{Nodes, topics and the two GPUs (desktop: registration;
laptop: tracking, refinement, planning)}{Software architecture of the cell.}

\section{Frames and Transformations}
\label{sec:sys_frames}
A part's pose in the robot base frame chains forward kinematics, the hand-eye extrinsic
and the camera's estimate:
\begin{equation}
\Tf{B}{O}(t) = \Tf{B}{E}\big(\vect{q}(t)\big)\;\Tf{E}{C}\;\Tf{C}{O}(t).
\label{eq:frame_chain}
\end{equation}
\TODO{Eye-in-hand consequences: everything stored in the base frame; TF looked up at
the image stamp; masks valid for one frame only.}
\SRC{README Algorithms \S1; admitans kontrol açıklama.md \S2.}
\figplaceholder{fig:frames}{TF tree: base\_link -- tool0 -- camera\_color\_optical\_frame --
object; the frame chain drawn on the robot}{Coordinate frames.}

\section{The Calibration Chain}
\label{sec:sys_calibration}
\TODO{The principle: measure each link against a reference that does not depend on it;
the robot (pen touches) is the reference. The order matters: a step on top of an
uncorrected earlier one learns its error.}
\begin{table}[htbp]
\caption{The calibration chain: each link, its reference and its residual.}
\label{tab:calib_chain}
\centering\small
\begin{tabular}{@{}llll@{}}
\toprule
Link & Method & Reference & Residual \\
\midrule
Kinematics & factory calibration in every tool & pendant TCP pose & $<0.1$ mm \\
Pen tip (TCP) & 4-point TCP; roll / tilt touches & table plane & 0.2 mm lateral \\
Table plane & 7 pen touches, plane fit & the robot & RMSE 0.22 mm \\
Camera depth & Tare with a robot-computed distance & table plane & 1.1 mm (0.3--0.65 m) \\
Extrinsic rotation & tilt split at 4 wrist yaws & table plane & $\approx 0.07^\circ$ \\
Extrinsic translation & online, from multi-view disagreement & views + table & \S\ref{sec:mv_selfcal} \\
\bottomrule
\end{tabular}
\end{table}

\subsection{Kinematics, Tool and Table}
\TODO{Factory calibration (nominal vs.\ calibrated: 2.4--4.2 mm at the tip); pen TCP;
table plane by touch (tilt 1.57$^\circ$).}
\SRC{README \S14 (29 Sep row), scripts/table\_touchoff.py, tcp\_offset\_probe.py.}

\subsection{Hand-Eye Calibration}
\TODO{AX = XB; the Park--Martin solve \cite{park1994handeye}; the residual judge; the
TCP-frame pitfall; the tilt separation at four yaws.}
\begin{equation}
\mathbf{A}_i\,\mathbf{X} = \mathbf{X}\,\mathbf{B}_i,\qquad \mathbf{X} = \Tf{E}{C}.
\label{eq:handeye}
\end{equation}
\SRC{README \S14 (24 Sep, 28 Sep rows); helper/calibration/; thesis\_notes.md.}

\subsection{Depth Calibration Against the Robot}
A disparity offset $\delta$ (pixels) in a stereo camera with focal length $f$ and baseline
$B$ produces a depth error that grows with the square of the range:
\begin{equation}
\Delta z = \frac{z^{2}\,\delta}{f\,B}.
\label{eq:depth_error}
\end{equation}
\TODO{Measurement at three heights; the Tare with a ground-truth distance computed from
the robot; the first Tare's 1.1\% scale error with a hand-measured distance.}
\SRC{README \S16; scripts/tare\_distance.py, extrinsic\_check.py; thesis\_notes.md
(calibration story).}
\figplaceholder{fig:depth_law}{Camera-vs-pen table offset at 0.29 / 0.44 / 0.63 m, before
and after the Tare, with the $a+kz^2$ fit}{Depth error against range, before and after
calibration.}
